\documentclass[10pt,a4paper]{article}
\usepackage[margin=18mm,headheight=15pt]{geometry}
\usepackage{amsmath,amssymb,booktabs,array,tabularx,xcolor,fancyhdr,draftwatermark,hyperref}
\usepackage[T1]{fontenc}\usepackage{lmodern}
\hypersetup{pdftitle={Experiment GUIDE v3: When the Method Should Win or Lose},pdfsubject={Conditional research guide, not observed results or fitted rankings},pdfauthor={Earning Roles project}}
\SetWatermarkText{GUIDE --- NOT OBSERVED}\SetWatermarkColor{red!14}
\SetWatermarkAngle{35}\SetWatermarkFontSize{32pt}\SetWatermarkScale{1}
\pagestyle{fancy}\fancyhf{}
\fancyhead[L]{\color{red!70!black}\bfseries GUIDE v3: CONDITIONS AND TRADEOFFS}
\fancyhead[R]{\color{red!70!black}NOT RESULTS}
\fancyfoot[L]{\color{red!70!black}2026-10-08 / UNFITTED HYPOTHESES / NOT FOR SUBMISSION}
\fancyfoot[R]{\thepage}
\setlength{\parindent}{0pt}\setlength{\parskip}{5pt}
\setlength{\tabcolsep}{4pt}\renewcommand{\arraystretch}{1.16}
\newcolumntype{L}[1]{>{\raggedright\arraybackslash}p{#1}}
\newcommand{\lead}[1]{\vspace{4pt}\textbf{#1}\par}
\newcommand{\rare}{\textsc{rare}}
\newcommand{\caution}{\textcolor{red!70!black}{\textbf{Directional hypotheses and algebraic cases; no fitted method ranking.}}\par}
\begin{document}
{\Large\bfseries When Should the Method Win---and Lose?}\par\caution
The objective is future team quality at complete cost, with reliable attribution
and timely updates. Reasonable expectations follow a mechanism and a matched
contrast. We do not assign a loss merely to make a results table look credible.
This guide supplements v2's exact conditional distribution model with conditions,
strong controls, adverse cases and falsifiers for the three body tables.

\lead{Independent review: the full forecast gate has not passed}
\begin{center}\small
\begin{tabularx}{\linewidth}{@{}L{.23\linewidth}L{.15\linewidth}X@{}}\toprule
Review dimension & Verdict & Evidence or missing requirement\\\midrule
Conditional arithmetic & Pass in scope & v2's same-outcome means, covariance and finite-count errors; assumed inputs.\\
Benchmark fit/authority & Partial / open & ArtifactRole fits the handoff question; distinct qualified roots and native later-use bridge remain open. PeerSelect is a separate adaptation track.\\
Baseline fairness & Open & Same information, task-dependent matching, eligibility, exploration and complete cost must be qualified live.\\
Empirical means/errors & Not identified & No qualified method-level future stream population or fitted uncertainty.\\
Advantage/weakness logic & Repaired in v3 & Two-sided conditions below; superiority still requires actual outcomes.\\\bottomrule
\end{tabularx}\end{center}

\lead{Existing evidence: snapshots, not a forecast sample}
The C1 snapshot cited by v2 has eight complete real API calls on one PIPE3 root,
no executed task after its updates, and a positive judgment despite failing
producer/terminal checks. Its 0.520/0.297 ms timings are single core updates.
The separate October 8 diagnosis made two real judgment calls, stopped on a
producer false acceptance, and left two cells unstarted. A public execution trace
exposes the datetime space/\texttt{T} mismatch; model use of that trace is untested.
Neither snapshot estimates judgment error rates or future method gains.

\lead{Table 1 baseline roles: a method must beat the right alternative}
\begin{center}\small
\begin{tabularx}{\linewidth}{@{}L{.27\linewidth}X@{}}\toprule
Body-table row / control & What the contrast must rule out\\\midrule
Uniform / no update & Exploration or base prior alone; cold start can tie with more overhead.\\
Raw acceptance / terminal only & Judgment content beyond raw verdicts or independent outcomes; neither is privileged truth.\\
Contextual trust linear & Ordinary contextual learning with the same public information; no presumed ranking.\\
Task-conditioned ridge$^{*}$ & Ordinary task--peer preference reversal; additive identity features alone are insufficient.\\
Pooled controller & Benefit of local personalized state versus pooling; sparse histories may favor pooling.\\
Meta-Team L2-public & Closest published public-state alternative; adapter remains unqualified.\\
\rare{} / matched composition$^{*}$ & Added mechanism beyond identical representation, gates and ordinary update/retrieval components.\\\bottomrule
\end{tabularx}\end{center}
\footnotesize $^{*}$Diagnostic comparators supplement the eight Table 1 rows;
this guide does not silently change the active paper matrix. Ridge has offline
algebraic controls, not live parity. Published native scores constrain metric
definitions and plausible scale, not our ranking: TeamBench partial $.7544$ is
not binary pass; Meta-Team's 13.1 pp system gain is not its 3.2 pp L2 increment.
The versioned reference archive preserves conflicting results.\normalsize

\clearpage
{\large\bfseries Table 1: Information and Future Allocation}\par\caution
Quality means an independently scored \emph{later} task. $U=Y-\lambda C$ uses
all declared costs. A changed assignment is a process diagnostic, not a quality
goal; frequent changes can be harmful. Every direction below is conditional.

\begin{center}\small
\begin{tabularx}{\linewidth}{@{}L{.20\linewidth}L{.22\linewidth}X L{.23\linewidth}@{}}\toprule
Condition & Matched contrast & Plausible gain and possible loss & Falsifying observation\\\midrule
Varied, informative judgments & Full vs count-only, J-masked, Qp-only and terminal-only & Lower held-out prediction loss and better later allocation \emph{if} J adds information. Extra judgment/retrieval cost may erase utility gain. & Masking J preserves the gain, or predictions improve but later quality does not.\\
Public execution facts explain the defect & Trace+J vs trace-only; J-masked & Added explanation may help transfer, but trace-only may tie or win at lower cost. No assumed J increment. & All improvement is reproduced by the same execution trace without J.\\
Stable task--peer crossover & Full vs task-conditioned ridge and matched composition & Context-specific matching may beat uniform/global trust. Advantage over the strong task-conditioned control is unknown. & Only weak additive controls lose, or ridge explains the whole gain.\\
Uninformative or constant J & Full vs no-update, count-only and strong prior & No J-specific gain is justified. If choices/outcomes are identical and extra cost is positive, full utility is lower. & Claimed J benefit persists after masking/permuting with valid support.\\
Incorrect, stale or mismatched evidence & Full vs static prior and same-info strong control & Learned evidence can reduce quality; costs and correction lag can worsen utility. A robust control may do better. & Benefit depends on wrong peer/version matching, leaked truth or future information.\\
Sparse comparable history / new root & Full vs pooled, prior and ridge; 0/1/K support & Personalized history may help once supported; pooling may win with few labels. Cold-start overhead and negative transfer are realistic risks. & Gain disappears under root/time holdout or only follows repeated task templates.\\\bottomrule
\end{tabularx}\end{center}

\lead{Fairness is part of the mechanism test}
Match the task, candidate@version/menu, public execution facts, feature/encoder
coverage, read cutoff, selected-only feedback, delay, eligibility/UNKNOWN,
propensity and exploration, capacity and complete service budget. Separate
two interventions: removing a responsibility gate tests that gate; changing the
updater while keeping the gate tests the updater. Do not give only full \rare{}
the execution trace or stronger task features.

\lead{Prediction comparisons cannot use different selected populations}
Compare probability forecasts on a common, sealed evaluation population with
legally acquired real labels and specified sampling weights/propensities.
Keep that evaluation separate from training feedback. Otherwise a policy may
change its own Brier by routing to easier cases. Choice propensity is not
predicted correctness. In v2, $p=q+e$ is a latent-oracle sensitivity case,
not a deployable predictor; S2/S3 Brier equality follows the symmetric $.7/.3$
assumption and is not a general routing property.

\clearpage
{\large\bfseries Tables 2--3: Adaptation, Safety and Mechanisms}\par\caution
\lead{Table 2: fast updates do not guarantee low end-to-end latency}
\begin{center}\small
\begin{tabularx}{\linewidth}{@{}L{.20\linewidth}L{.23\linewidth}X@{}}\toprule
Condition & Comparator & Two-sided expectation / required measurement\\\midrule
Stable peers, little novelty & No history / static prior / ordinary updater & Full may tie quality and lose on service/state cost. Measure whole-service p50/p95, state and complete cost; core time alone is insufficient.\\
Gradual drift, timely legal labels & Frozen prior / same-info RLS, SGD or refit & Online feedback may recover earlier than frozen state. A strong updater may match or beat full; retention can slow adaptation. Measure old/new holdouts, recovery and utility.\\
Abrupt change / peer replacement & Static prior / reset / version-aware control & Stale profiles may cause negative transfer; version isolation may avoid it but restart learning. Report recovery, old-task loss and feedback coverage.\\
Long feedback delay / few eligible outcomes & History-only / matched delayed updater & Public observations may help early decisions if predictive; persistent training still waits for legal later credit. Sparse credits can make adaptation slower.\\
High arrival rate / expensive evidence & Local uniform / matched service stack & Selection quality can rise while backlog and p95 worsen. Measure arrivals, full service distribution, queueing, drops/timeouts and net payoff.\\\bottomrule
\end{tabularx}\end{center}
Graph-IPD is a native multi-round cooperation reference, not an ArtifactRole
score or an online service result. Its published cooperation bounds cannot
populate payoff, update p95, forgetting or recovery cells. Timing, drift and
retention methods are not finally locked; these are tests, not feature claims.

\lead{Table 3: attribution safety and organic interaction}
\begin{center}\small
\begin{tabularx}{\linewidth}{@{}L{.23\linewidth}X L{.25\linewidth}@{}}\toprule
Contrast & What can improve / what can worsen & Disconfirming outcome\\\midrule
Gate vs no gate & Independent owner truth can reduce false producer credit. Abstention costs coverage, usable labels and update speed. Legal positive/negative J remain observations; gate is not accept-only. & Same safety but less coverage and worse future utility; factually wrong J still treated as reward.\\
Public evidence $\times$ delayed credit (four cells) & Early public information plus later correction may work together. Neither interaction nor dominance of the full cell is automatic. & No joint increment, or full loses to one component or matched composition.\\
Local vs nonlocal & Locality can lower communication/cost but miss a better distant peer. Compare complete costs and quality; no universally best neighborhood. & Cost saving comes entirely from a smaller information/compute budget.\\
Correction/replay vs ordinary state & Exact identity/idempotence can prevent duplicate credit; retention/retraction has state and service costs. Engineering safety is not an efficacy estimate. & Legal credit counted twice, correction changes the wrong version, or gains vanish with matched replay.\\\bottomrule
\end{tabularx}\end{center}
Four-cell utility interaction is $I_U=U_{11}-U_{10}-U_{01}+U_{00}$, with
$\operatorname{Var}(\widehat I_U)=c^\top\Sigma c$ for $c=(1,-1,-1,1)$ in that
cell order. Report the qualified stream covariance; do not assign four convenient
SDs. Zero interaction does not by itself negate the value of components, but
fails a claimed synergy. All adaptive arms keep their own actual histories.

\clearpage
{\large\bfseries Quantitative Boundaries and Distribution Acceptance}\par\caution
\lead{A quality gain can lose after complete cost}
\[
 \Delta U=\Delta\mathbb E[Y]-\lambda\Delta\mathbb E[C],\qquad
 \Delta U>0\ \Longleftrightarrow\ \Delta Q>\lambda\Delta C.
\]
Case: $\lambda=.1$ below is the v2 \emph{illustrative} weight, not a frozen
scientific acceptance margin. Cost is normalized complete cost. These are
algebraic boundary inputs, not fitted means, predictions or method ranks.
\begin{center}\small
\begin{tabularx}{\linewidth}{@{}lXrrr@{}}\toprule
Case & Explicit assumption & $\Delta Q$ & $\Delta C$ & $\Delta U$\\\midrule
@@BOUNDARY_ROWS@@
\bottomrule\end{tabularx}\end{center}
No case is assigned to \rare{} or a baseline. Observed methods may cross these
boundaries; that is exactly what the complete-cost experiment must determine.

\lead{Why v2's repair-only cost law is optimistic for routing}
Within v2 only, $\Delta U=(1+\lambda\gamma/c_0)d\Delta\kappa
-(\lambda/c_0)\Delta h$. With $d=.2$, $\gamma=1$, $c_0=1.5$ and $\lambda=.1$,
positive utility requires $\Delta\kappa>@@KAPPA_THRESHOLD@@\Delta h$.
Case: the assumed S3--S2 routing difference $.25$ reverses only when extra raw
overhead exceeds $@@H_THRESHOLD@@$. This weak penalty is not empirically justified.
Real selection, judging, retrieval, publication, repair, failures and resource
use must enter the cost ledger; stronger peers may be more expensive even when
successful. Refit the model rather than forcing reality into its negative
quality--cost correlation.

\lead{UNKNOWN can reverse an apparent quality ranking}
For all assignments, let $m$ be the fraction with unknown binary outcome and
$q_{\rm obs}$ the success rate among known outcomes. Without a missingness model,
\[
 Q\in[(1-m)q_{\rm obs},\ (1-m)q_{\rm obs}+m].
\]
Case: hypothetical candidate $.65$ with $m=.15$ gives
$[@@UNKNOWN_CANDIDATE@@]$; control $.60$ with $m=.02$ gives
$[@@UNKNOWN_CONTROL@@]$. These overlap, so the higher observed rate does not
establish a robust advantage. The bounds are logical limits, not confidence
intervals; sampling uncertainty still applies. Unscorable resource failures
remain UNKNOWN, with their attempted costs and denominators retained.

\lead{Means, errors and distributions must be validated together}
\small
1. Seal qualified roots, independent live histories, same-info comparators,
primary quality/cost metrics, budgets and stopping rules before a pilot.
Fit on development only; validate on untouched future streams/root holdouts.\par
2. Model the joint observable outcomes and complete costs. Utility variance is
$\sigma_U^2=\sigma_Q^2+\lambda^2\sigma_C^2-2\lambda\operatorname{Cov}(Q,C)$.
High quality--utility correlation is not independent corroboration. No guessed
per-cell SD or 95\% interval is supplied here.\par
3. Report residual bias, root heterogeneity, serial correlation, covariance,
80/95\% prediction coverage \emph{and} width with uncertainty, alongside
UNKNOWN by arm/root/J/owner. Prediction intervals differ from confidence
intervals of a treatment effect; small root counts cannot justify a normal
1.96 band automatically.\par
4. Test the favorable and adverse conditions above. A strong control explaining
the whole gain, or complete cost removing it, triggers diagnosis and method
iteration. It does not authorize rewriting Goal or discarding failure logs.\normalsize

\end{document}
